\declaredin{BPatch\_type.h}

This class represents a type in the mutatee. Anything that reports the type
of something returns one, such as \code{getReturnType} on a
\code{BPatch\_function} or \code{getType} on a \code{BPatch\_variableExpr}, and
the \code{BPatch::create*} factory methods build them. The rules Dyninst uses
to decide whether two types are compatible are given in
Section~\ref{sec:typesystem}.

\begin{apient}
const char* getName() const;
\end{apient}

\apidesc{
  Return the name of the type.
}

\begin{apient}
BPatch_dataClass getDataClass() const;
\end{apient}

\apidesc{
  Return the data class of the type, which is what decides how it may be
  used and which of the methods below apply to it. The values are listed in
  Section~\ref{sec:typesystem}.
}

\begin{apient}
unsigned int getSize();
\end{apient}

\apidesc{
  Return the size of the type in bytes.
}

\begin{apient}
int getID() const;
\end{apient}

\apidesc{
  Return the identifier Dyninst gave this type. Identifiers are unique within
  a program. To reach other types, use \code{getConstituentType} and
  \code{getComponents} rather than looking one up by its identifier.
}

\begin{apient}
virtual bool operator==(const BPatch_type& otype) const;
\end{apient}

\apidesc{
  Return \code{true} if the two are the same type: same identifier, same
  data class, and the same underlying \code{SymtabAPI} type. This is a
  stricter test than \code{isCompatible}, which asks only whether a value of
  one may be used where the other is wanted.
}

\begin{apient}
bool isCompatible(BPatch_type* otype);
\end{apient}

\apidesc{
  Return \code{true} if this type is compatible with \code{otype} under the
  rules in Section~\ref{sec:typesystem}.
}

\begin{apient}
unsigned long getLow() const;
unsigned long getHigh() const;
\end{apient}

\apidesc{
  Return the lower and upper bounds of a ranged type, such as an array or a
  subrange. Both return 0 for a type that has no range, so a 0 does not
  distinguish a bound of zero from a type these do not apply to; use
  \code{getDataClass} to tell the two apart.
}

\begin{apient}
BPatch_type* getConstituentType() const;
\end{apient}

\apidesc{
  Return the type this one is built from: the target of a pointer or
  reference, the type a typedef names, the underlying type of an enumeration,
  or the element type of an array. Returns \code{nullptr} for a type that is
  not built from another.
}

\begin{apient}
std::vector<BPatch_field*>* getComponents() const;
\end{apient}

\apidesc{
  Return the fields of a structure, union or common block. For an enumerated
  type the identifiers are returned as fields, each carrying its value. For a
  pointer, reference or typedef the fields of the constituent type are
  returned. Returns \code{nullptr} for any other type. The vector is newly
  allocated for the caller.
}

\begin{apient}
Dyninst::SymtabAPI::Type* getSymtabType() const;
\end{apient}

\apidesc{
  Return the \code{SymtabAPI} type this one wraps, which carries more of the
  type's detail than \code{BPatch\_type} exposes. \code{SymtabAPI::convert},
  listed in Section~\ref{sec:conversions}, forwards to this method, so the two
  return the same pointer.
}

\begin{apient}
std::vector<BPatch_cblock*>* getCblocks() const;
\end{apient}

\apidesc{
  Intended to return the common block objects for a common block type, the
  way \code{typeCommon} does in SymtabAPI. It does not. For anything that is
  not a common block type it returns \code{nullptr}, and for one that is it
  returns a vector it never puts anything into, so the result is always empty.
  Until that is fixed, reach the common blocks through \code{getSymtabType}
  and ask \code{typeCommon} for them.
}
